% Generated draft: requires mathematical and source review.
\documentclass[11pt]{article}
\usepackage[margin=25mm]{geometry}
\usepackage{fontspec}
\setmainfont{Latin Modern Roman}
\usepackage{amsmath,amssymb,mathtools,unicode-math}
\setmathfont{Latin Modern Math}
\usepackage{enumitem,needspace,xurl,hyperref}
\setlength{\parindent}{0pt}
\setlength{\parskip}{5pt}
\setlength{\emergencystretch}{3em}
\newcommand{\N}{\mathbb{N}}
\newcommand{\Z}{\mathbb{Z}}
\newcommand{\Q}{\mathbb{Q}}
\newcommand{\R}{\mathbb{R}}
\newcommand{\C}{\mathbb{C}}
\newcommand{\F}{\mathbb{F}}
\newcommand{\A}{\mathbb{A}}
\newcommand{\PP}{\mathbb{P}}
\newcommand{\E}{\mathbb{E}}
\newcommand{\eps}{\varepsilon}
\newcommand{\dd}{\,\mathrm{d}}
\begin{document}
\section*{UnsolvedMath ID: 9400112}
\textbf{Source problem code:} catalogue entry
\subsection{Definitions and mathematical statement}
A perfect integer cube is $x^3$ with $x\in\mathbb Z$; negative cubes and zero are allowed. Is it true that
\[
\forall n\in\mathbb Z\ \exists x_1,x_2,x_3,x_4\in\mathbb Z,
\qquad n=x_1^3+x_2^3+x_3^3+x_4^3\,?
\]
The summands need not be distinct. With zero permitted, exactly four slots are equivalent to at most four cubes. A restriction to strictly positive cubes would be a different statement and is not adopted.

The supplied metadata labels this solved and its research summary attributes the signed four-cube assertion to Linnik, but supplies no primary proof citation. The same summary discusses seven- and nine-positive-cube statements, which are different quantified problems and are not provided as a proof of the four-signed-cube target. The pinned Wikipedia provenance instead presents this explicit bullet as unsolved. These conflicting/unverified status claims are retained without asserting a resolution.
\subsection{Short English statement}
Can every integer be expressed as four cubes of integers, allowing negative values, zero, and repeated summands? The source's solved attribution has no primary proof citation and conflicts with its unsolved-list provenance.
\subsection{Sources}
\begin{itemize}
\item Wikipedia contributors, \emph{List of unsolved problems in mathematics}, Diophantine equations bullet 112, wikitext line 141, pinned revision 1366636767, \url{https://en.wikipedia.org/w/index.php?title=List_of_unsolved_problems_in_mathematics&oldid=1366636767#Number_theory}; accessed 2026-07-29, MediaWiki source extraction.
\item The supplied summary attributes a four-signed-cube result to Linnik and refers to Wikipedia pages ``Waring's problem,'' ``Linnik's theorem,'' and ``Sum of cubes,'' without a primary publication citation or theorem statement supporting that attribution.
\item Local registry \texttt{data/unsolvedmath/problems.json}, permanent UnsolvedMath ID $9400112$, source code catalogue entry, \emph{Can every integer be written as a sum of four perfect cubes}. Dataset attribution: Ulam AI and the UnsolvedMath Contributors.
\end{itemize}
\end{document}
